% GENERATED FILE. Edit canonical lessons, then run npm run generate:books.
% canonical-source-hash: fnv1a64:70fcd6b953ca85c7
% canonical-lessons: ML-C18-oru, ML-C18-mani, ML-S121-letter-ssa

\chapter{Telling Time}
\label{ch:ml-telling-time}
\hlchaptermodality{18}

\begin{chapteropening}
\textbf{By the end of this chapter:} \emph{I can say the hour in Malayalam with maṇi, as in oru maṇi, one o'clock.}

Say maṇi for hour, bell and gem, then oru maṇi, and state where Malayalam's and Tamil's dictionaries disagree about it.
\end{chapteropening}

\section[\texorpdfstring{\ml{ഒരു} (oru)}{ഒരു (oru)}]{\ml{ഒരു} — one, when something follows it}

\label{lesson:ML-C18-oru}



\begin{quote}
Count one to five out loud.

Now say the first of them on its own: \textbf{\ml{ഒന്ന്}}.
\end{quote}

\subsection*{What to know first}
\textbf{\ml{ഒരു}} (\emph{oru}) — \textbf{one}, with a noun coming after it.

\textbf{\ml{ഒന്ന്}} and \textbf{\ml{ഒരു}} are one number wearing two shapes, and which shape you reach for depends on nothing more than whether a noun follows it:

\noindent
\begin{tabularx}{\linewidth}{@{}>{\raggedright\arraybackslash}X>{\raggedright\arraybackslash}X@{}}
\toprule
\textbf{} & \textbf{} \\
\midrule
\textbf{\ml{ഒന്ന്}} & standing alone — \emph{how many? one} \\
\textbf{\ml{ഒരു}} & leaning on a noun — \textbf{\ml{ഒരു} \ml{പേര്}}, \emph{one name} \\
\bottomrule
\end{tabularx}

You counted with the first. The second is the one you will meet on nearly every page from here, because a number usually has something to count.

\textbf{Tamil keeps the same pair}, \emph{oru} against \emph{oṉṟu}, split the same way — which is a sign the split is old rather than a tidy-up Malayalam did on its own.

\begin{cousinweb}
\textbf{\ml{ഒരു}} (\emph{oru}) is native Dravidian, the shape \textquotedblleft{}one\textquotedblright{} takes in front of a noun. DEDR 990 lists Tamil \emph{oru} before a consonant and \emph{ōr} before a vowel, with Kannada \emph{ondu}, \textquotedblleft{}one thing\textquotedblright{}.
\end{cousinweb}

\begin{grammarlens}[title={Grammar Lens: the same word does the work of \textquotedblleft{}a\textquotedblright{}}]
\textbf{Malayalam has no word for \emph{a}.} A bare noun covers it: \textbf{\ml{പേര്}} on its own is \emph{name}, \emph{a name}, or \emph{the name}, and which one it is comes from what is being talked about.

But when a speaker wants to set \textbf{one unspecified thing} in front of you — the \emph{a} that means \emph{some one or other} — \textbf{\ml{ഒരു}} is what they say. So the same small word carries two weights:

\noindent
\begin{tabularx}{\linewidth}{@{}>{\raggedright\arraybackslash}X>{\raggedright\arraybackslash}X@{}}
\toprule
\textbf{} & \textbf{} \\
\midrule
counting & \textbf{\ml{ഒരു} \ml{പേര്}} — \emph{one name} \\
introducing & \textbf{\ml{ഒരു} \ml{പേര്}} — \emph{a name} \\
\bottomrule
\end{tabularx}

\textbf{Nothing in the word changes between those two rows.} The difference lives in why it was said, and a listener takes it from the situation — the same bargain Malayalam makes everywhere else it declines to mark something.

\textbf{The next lesson puts a noun after it straight away}, and when it does, this is the shape you will be reading.
\end{grammarlens}

\subsection*{Your turn}
Say these aloud:

\begin{itemize}
\raggedright
  \item \textbf{\ml{ഒന്ന്}}, then \textbf{\ml{ഒരു}} — and say which of the two wants a noun after it
  \item \textbf{\ml{ഒരു} \ml{പേര്}}, and both things it could mean
  \item what a bare \textbf{\ml{പേര്}} could mean with nothing in front of it
\end{itemize}

\begin{tcolorbox}[breakable,colback=teal!4,colframe=teal!35!black,title={Before you move on}]
Which shape of \emph{one} stands alone? (\textbf{\ml{ഒന്ന്}}.) Which leans on a noun? (\textbf{\ml{ഒരു}}.) Does Malayalam have a word for \emph{a}? (\textbf{No — but \ml{ഒരു} does that work when a speaker wants it done.})
\end{tcolorbox}
\section[\texorpdfstring{\ml{മണി} (maṇi)}{മണി (maṇi)}]{\ml{മണി} — the bell that tells the hour}

\label{lesson:ML-C18-mani}



\begin{quote}
An hour you can hear: before clocks were everywhere, a bell struck it. Malayalam still names the hour after the bell.
\end{quote}

\subsection*{What to know first}
\textbf{\ml{മണി}} (\emph{maṇi}) = \textquotedblleft{}\textbf{hour, o'clock}\textquotedblright{}, and also \textquotedblleft{}\textbf{bell, gong}\textquotedblright{}. The bell came first: the hour is the stroke of the bell that announces it.

\begin{culture}
The Dravidian Etymological Dictionary lists this word as Dravidian (DEDR 4672): Tamil \emph{maṇi}, \textquotedblleft{}bell, gong, hour\textquotedblright{}, Malayalam \emph{maṇi}, \textquotedblleft{}little bells, gong, hour by the bell\textquotedblright{}, with Kota, Toda and Tulu bell words beside them. DEDR then adds a query: perhaps from Sanskrit \emph{maṇi}, \textquotedblleft{}jewel\textquotedblright{}. It does not decide. Malayalam's own row hints at how the two could meet, since its little bells were worn as jewels, but that link stays an open question, not a finding.
\end{culture}

\begin{grammarlens}[title={Grammar Lens: telling the time}]
\begin{quote}
\textbf{\ml{ഒരു} \ml{മണി}.} — \textquotedblleft{}One o'clock.\textquotedblright{} (\emph{oru maṇi}, \textquotedblleft{}one hour/bell\textquotedblright{}) \textbf{\ml{രണ്ട്} \ml{മണി}.} — \textquotedblleft{}Two o'clock.\textquotedblright{} (\emph{raṇṭŭ maṇi}, \textquotedblleft{}two hour/bell\textquotedblright{})
\end{quote}
\end{grammarlens}

\subsection*{Your turn}
Say these aloud:

\begin{itemize}
\raggedright
  \item \textquotedblleft{}maṇi\textquotedblright{} — hour, and the bell that strikes it
  \item \textquotedblleft{}oru maṇi\textquotedblright{} — one o'clock
  \item the honest point — DEDR lists maṇi as a Dravidian bell word shared with Tamil, and only queries a Sanskrit \textquotedblleft{}jewel\textquotedblright{} source
\end{itemize}

\begin{tcolorbox}[breakable,colback=teal!4,colframe=teal!35!black,title={Before you move on}]
What did \emph{maṇi} name before it named the hour? (\textbf{A bell or gong}.) Which language family does DEDR put it in, and with which Tamil word? (\textbf{Dravidian}, with Tamil \emph{maṇi}, \textquotedblleft{}bell, gong, hour\textquotedblright{}.) Is a link to Sanskrit \emph{maṇi}, \textquotedblleft{}jewel\textquotedblright{}, settled? (\textbf{No} — DEDR marks it with a query.)
\end{tcolorbox}
\section[\texorpdfstring{\ml{ഷ} (ṣa)}{ഷ (ṣa)}]{\ml{ഷ} — one character, met inside words you already say}

\label{lesson:ML-S121-letter-ssa}



\begin{quote}
Before the new one: \ml{◌ം} — what does it do?

One character this time. Just one — and you have been saying it for pages without knowing which mark on the page it was.
\end{quote}

\subsection*{Script you'll notice: \ml{ഷ}}
\textbf{\ml{ഷ}} — \emph{ṣa}.

It is a \textbf{consonant}, and in this script a consonant is never bare: it comes with an \emph{a} already in it. So it is not \emph{ṣ}, it is \textbf{ṣa}.

You already say these, and every one of them has \ml{ഷ} somewhere inside it:

\begin{itemize}
\raggedright
  \item \textbf{\ml{സന്തോഷം}} \emph{santōṣaṁ} — joy / pleased to meet you
  \item \textbf{\ml{ക്ഷമിക്കണം}} \emph{kṣamikkaṇaṁ} — (you) should forgive / sorry
\end{itemize}

\subsection*{Writing: \ml{ഷ} — follow the numbered strip}
\hlblockfigure{\detokenize{figures/ML-S121-letter-ssa-filmstrip.pdf}}{How \ml{ഷ} is written, stroke by stroke}

Follow the numbered strip: start where movement 1 begins, and let each panel show you the next movement before your pen makes it. Then write \ml{ഷ} yourself — slowly, and larger than it is printed.

\begin{quote}
The strip shows \textbf{where to start the character and which way to travel}, in one attested order. That is taught with real variation from school to school, so the strip names its source beneath it: treat it as a sound way in, not the only one.
\end{quote}

\subsection*{Your turn}
\begin{itemize}
\raggedright
  \item \emph{Look:} at these words, and find \ml{ഷ} in the ones that have it
\end{itemize}

\begin{quote}
\ml{സന്തോഷം}  ·  \ml{ക്ഷമിക്കണം}  ·  \ml{നമസ്കാരം}
\end{quote}

\begin{itemize}
\raggedright
  \item \emph{Trace it:} \ml{ഷ} three times, saying \emph{ṣa} as you finish each one
  \item \emph{Look:} back at any page of this chapter and find \ml{ഷ} once more
\end{itemize}

\begin{tcolorbox}[breakable,colback=teal!4,colframe=teal!35!black,title={Before you move on}]
Which character is this — \ml{ഷ}? What sound does it carry? (\emph{\textbf{ṣa}}.) Name one word you already say that contains it.
\end{tcolorbox}
